\section{晚清改革其余账本事件锚点}
本节补齐晚清改革期尚未导航的账本条目，区分制度公布、资金与人事、地方采纳、公共信息和社会结果。
\subsection{1870—1879}
\eventanchor{EQ-1875-YUNNAN-MUSLIM-WAR-PRELUDE}{「杜文秀政权覆亡前后，云南回…}光绪元年（1875年），「杜文秀政权覆亡前后，云南回民战争进入惨烈收束阶段」之前的条件、信息和争议已通过中央与地方材料显现，构成该事件可追踪的前史节点。核心取证：《清德宗实录》光绪元年相关条；军机处档、战事方略及地方奏报；同时期地方档案、地方志、日记或对方材料应按具体地区补检。该项锚定的是前史条件与信息背景，不把条件直接写成唯一因果。来源保留：官修编年和地方材料的记录密度与立场并不相同；本行不将条件、传闻、呈报或后见解释直接等同为确定因果。
\eventanchor{EQ-1875-YUNNAN-MUSLIM-WAR-DECISION}{围绕「杜文秀政权覆亡前后，云…}光绪元年（1875年），围绕「杜文秀政权覆亡前后，云南回民战争进入惨烈收束阶段」的命令、调度、照会或呈报形成独立的中央—地方行动文书链。核心取证：《清德宗实录》光绪元年相关条；军机处档、战事方略及地方奏报。该项锚定的是命令或决策环节，命令抵达和结果仍须另外核验。来源保留：奏折、上谕、部议、军机档或条约可证明行政动作和机构立场，不自动证明所有地区、群体或对手按其意图行动。
\eventanchor{EQ-1875-YUNNAN-MUSLIM-WAR-AFTERMATH}{「杜文秀政权覆亡前后，云南回…}光绪元年（1875年），「杜文秀政权覆亡前后，云南回民战争进入惨烈收束阶段」引发的善后、执行或地方回应构成可由不同材料追踪的后续节点。核心取证：《清德宗实录》光绪元年相关条；军机处档、战事方略及地方奏报；相关地区的档案、志书、契约、报刊或对方材料。该项锚定的是后续追踪问题，影响范围须按地区与群体分别判断。来源保留：战后、改革后或条约后的记忆、地方记录和官方总结常具有事后选择性；后续影响须按地点、群体和时间分层，不作整体化推断。
\eventanchor{EQ-1876-CHEFOO-CONVENTION-PRELUDE}{「《烟台条约》签订，扩大英国…}光绪二年（1876年），「《烟台条约》签订，扩大英国在华通商和内地活动权利」之前的条件、信息和争议已通过中央与地方材料显现，构成该事件可追踪的前史节点。核心取证：《清德宗实录》光绪二年相关条；《筹办夷务始末》相关年卷与中外照会、条约文本；同时期地方档案、地方志、日记或对方材料应按具体地区补检。该项锚定的是前史条件与信息背景，不把条件直接写成唯一因果。来源保留：官修编年和地方材料的记录密度与立场并不相同；本行不将条件、传闻、呈报或后见解释直接等同为确定因果。
\eventanchor{EQ-1876-CHEFOO-CONVENTION-DECISION}{围绕「《烟台条约》签订，扩大…}光绪二年（1876年），围绕「《烟台条约》签订，扩大英国在华通商和内地活动权利」的命令、调度、照会或呈报形成独立的中央—地方行动文书链。核心取证：《清德宗实录》光绪二年相关条；《筹办夷务始末》相关年卷与中外照会、条约文本。该项锚定的是命令或决策环节，命令抵达和结果仍须另外核验。来源保留：奏折、上谕、部议、军机档或条约可证明行政动作和机构立场，不自动证明所有地区、群体或对手按其意图行动。
\eventanchor{EQ-1876-CHEFOO-CONVENTION-AFTERMATH}{「《烟台条约》签订，扩大英国…}光绪二年（1876年），「《烟台条约》签订，扩大英国在华通商和内地活动权利」引发的善后、执行或地方回应构成可由不同材料追踪的后续节点。核心取证：《清德宗实录》光绪二年相关条；《筹办夷务始末》相关年卷与中外照会、条约文本；相关地区的档案、志书、契约、报刊或对方材料。该项锚定的是后续追踪问题，影响范围须按地区与群体分别判断。来源保留：战后、改革后或条约后的记忆、地方记录和官方总结常具有事后选择性；后续影响须按地点、群体和时间分层，不作整体化推断。
\eventanchor{EQ-1876-NORTH-CHINA-FAMINE-PRELUDE}{「丁戊奇荒在华北及西北造成广…}光绪二年（1876年），「丁戊奇荒在华北及西北造成广泛饥荒、迁徙和救济行动」之前的条件、信息和争议已通过中央与地方材料显现，构成该事件可追踪的前史节点。核心取证：《清德宗实录》光绪二年相关条；地方志、督抚奏折及地方档案；同时期地方档案、地方志、日记或对方材料应按具体地区补检。该项锚定的是前史条件与信息背景，不把条件直接写成唯一因果。来源保留：官修编年和地方材料的记录密度与立场并不相同；本行不将条件、传闻、呈报或后见解释直接等同为确定因果。
\eventanchor{EQ-1876-NORTH-CHINA-FAMINE-DECISION}{围绕「丁戊奇荒在华北及西北造…}光绪二年（1876年），围绕「丁戊奇荒在华北及西北造成广泛饥荒、迁徙和救济行动」的命令、调度、照会或呈报形成独立的中央—地方行动文书链。核心取证：《清德宗实录》光绪二年相关条；地方志、督抚奏折及地方档案。该项锚定的是命令或决策环节，命令抵达和结果仍须另外核验。来源保留：奏折、上谕、部议、军机档或条约可证明行政动作和机构立场，不自动证明所有地区、群体或对手按其意图行动。
\eventanchor{EQ-1876-NORTH-CHINA-FAMINE-AFTERMATH}{「丁戊奇荒在华北及西北造成广…}光绪二年（1876年），「丁戊奇荒在华北及西北造成广泛饥荒、迁徙和救济行动」引发的善后、执行或地方回应构成可由不同材料追踪的后续节点。核心取证：《清德宗实录》光绪二年相关条；地方志、督抚奏折及地方档案；相关地区的档案、志书、契约、报刊或对方材料。该项锚定的是后续追踪问题，影响范围须按地区与群体分别判断。来源保留：战后、改革后或条约后的记忆、地方记录和官方总结常具有事后选择性；后续影响须按地点、群体和时间分层，不作整体化推断。
\eventanchor{EQ-1878-Zuo-ZONGTANG-XINJIANG-PRELUDE}{「左宗棠军收复新疆大部，阿古…}光绪四年（1878年），「左宗棠军收复新疆大部，阿古柏政权崩解」之前的条件、信息和争议已通过中央与地方材料显现，构成该事件可追踪的前史节点。核心取证：《清德宗实录》光绪四年相关条；军机处档、战事方略及地方奏报；同时期地方档案、地方志、日记或对方材料应按具体地区补检。该项锚定的是前史条件与信息背景，不把条件直接写成唯一因果。来源保留：官修编年和地方材料的记录密度与立场并不相同；本行不将条件、传闻、呈报或后见解释直接等同为确定因果。
\eventanchor{EQ-1878-Zuo-ZONGTANG-XINJIANG-DECISION}{围绕「左宗棠军收复新疆大部，…}光绪四年（1878年），围绕「左宗棠军收复新疆大部，阿古柏政权崩解」的命令、调度、照会或呈报形成独立的中央—地方行动文书链。核心取证：《清德宗实录》光绪四年相关条；军机处档、战事方略及地方奏报。该项锚定的是命令或决策环节，命令抵达和结果仍须另外核验。来源保留：奏折、上谕、部议、军机档或条约可证明行政动作和机构立场，不自动证明所有地区、群体或对手按其意图行动。
\eventanchor{EQ-1878-Zuo-ZONGTANG-XINJIANG-AFTERMATH}{「左宗棠军收复新疆大部，阿古…}光绪四年（1878年），「左宗棠军收复新疆大部，阿古柏政权崩解」引发的善后、执行或地方回应构成可由不同材料追踪的后续节点。核心取证：《清德宗实录》光绪四年相关条；军机处档、战事方略及地方奏报；相关地区的档案、志书、契约、报刊或对方材料。该项锚定的是后续追踪问题，影响范围须按地区与群体分别判断。来源保留：战后、改革后或条约后的记忆、地方记录和官方总结常具有事后选择性；后续影响须按地点、群体和时间分层，不作整体化推断。
\eventanchor{EQ-1878-XINJIANG-PROVINCE-DEBATE-PRELUDE}{「朝廷讨论新疆设省、军政与财…}光绪四年（1878年），「朝廷讨论新疆设省、军政与财政安排」之前的条件、信息和争议已通过中央与地方材料显现，构成该事件可追踪的前史节点。核心取证：《清德宗实录》光绪四年相关条；宫中朱批奏折、内阁大库题本与《清史稿》相关志传；同时期地方档案、地方志、日记或对方材料应按具体地区补检。该项锚定的是前史条件与信息背景，不把条件直接写成唯一因果。来源保留：官修编年和地方材料的记录密度与立场并不相同；本行不将条件、传闻、呈报或后见解释直接等同为确定因果。
\eventanchor{EQ-1878-XINJIANG-PROVINCE-DEBATE-DECISION}{围绕「朝廷讨论新疆设省、军政…}光绪四年（1878年），围绕「朝廷讨论新疆设省、军政与财政安排」的命令、调度、照会或呈报形成独立的中央—地方行动文书链。核心取证：《清德宗实录》光绪四年相关条；宫中朱批奏折、内阁大库题本与《清史稿》相关志传。该项锚定的是命令或决策环节，命令抵达和结果仍须另外核验。来源保留：奏折、上谕、部议、军机档或条约可证明行政动作和机构立场，不自动证明所有地区、群体或对手按其意图行动。
\eventanchor{EQ-1878-XINJIANG-PROVINCE-DEBATE-AFTERMATH}{「朝廷讨论新疆设省、军政与财…}光绪四年（1878年），「朝廷讨论新疆设省、军政与财政安排」引发的善后、执行或地方回应构成可由不同材料追踪的后续节点。核心取证：《清德宗实录》光绪四年相关条；宫中朱批奏折、内阁大库题本与《清史稿》相关志传；相关地区的档案、志书、契约、报刊或对方材料。该项锚定的是后续追踪问题，影响范围须按地区与群体分别判断。来源保留：战后、改革后或条约后的记忆、地方记录和官方总结常具有事后选择性；后续影响须按地点、群体和时间分层，不作整体化推断。
\eventanchor{EQ-1879-RYUKYU-ANNEXATION-PRELUDE}{「日本吞并琉球，清廷的交涉未…}光绪五年（1879年），「日本吞并琉球，清廷的交涉未能恢复旧有册封关系」之前的条件、信息和争议已通过中央与地方材料显现，构成该事件可追踪的前史节点。核心取证：《清德宗实录》光绪五年相关条；《筹办夷务始末》相关年卷与中外照会、条约文本；同时期地方档案、地方志、日记或对方材料应按具体地区补检。该项锚定的是前史条件与信息背景，不把条件直接写成唯一因果。来源保留：官修编年和地方材料的记录密度与立场并不相同；本行不将条件、传闻、呈报或后见解释直接等同为确定因果。
\eventanchor{EQ-1879-RYUKYU-ANNEXATION-DECISION}{围绕「日本吞并琉球，清廷的交…}光绪五年（1879年），围绕「日本吞并琉球，清廷的交涉未能恢复旧有册封关系」的命令、调度、照会或呈报形成独立的中央—地方行动文书链。核心取证：《清德宗实录》光绪五年相关条；《筹办夷务始末》相关年卷与中外照会、条约文本。该项锚定的是命令或决策环节，命令抵达和结果仍须另外核验。来源保留：奏折、上谕、部议、军机档或条约可证明行政动作和机构立场，不自动证明所有地区、群体或对手按其意图行动。
\eventanchor{EQ-1879-RYUKYU-ANNEXATION-AFTERMATH}{「日本吞并琉球，清廷的交涉未…}光绪五年（1879年），「日本吞并琉球，清廷的交涉未能恢复旧有册封关系」引发的善后、执行或地方回应构成可由不同材料追踪的后续节点。核心取证：《清德宗实录》光绪五年相关条；《筹办夷务始末》相关年卷与中外照会、条约文本；相关地区的档案、志书、契约、报刊或对方材料。该项锚定的是后续追踪问题，影响范围须按地区与群体分别判断。来源保留：战后、改革后或条约后的记忆、地方记录和官方总结常具有事后选择性；后续影响须按地点、群体和时间分层，不作整体化推断。
\subsection{1880—1889}
\eventanchor{EQ-1880-ILI-NEGOTIATIONS-PRELUDE}{「清俄围绕伊犁归还和边界条件…}光绪六年（1880年），「清俄围绕伊犁归还和边界条件持续谈判」之前的条件、信息和争议已通过中央与地方材料显现，构成该事件可追踪的前史节点。核心取证：《清德宗实录》光绪六年相关条；《筹办夷务始末》相关年卷与中外照会、条约文本；同时期地方档案、地方志、日记或对方材料应按具体地区补检。该项锚定的是前史条件与信息背景，不把条件直接写成唯一因果。来源保留：官修编年和地方材料的记录密度与立场并不相同；本行不将条件、传闻、呈报或后见解释直接等同为确定因果。
\eventanchor{EQ-1880-ILI-NEGOTIATIONS-DECISION}{围绕「清俄围绕伊犁归还和边界…}光绪六年（1880年），围绕「清俄围绕伊犁归还和边界条件持续谈判」的命令、调度、照会或呈报形成独立的中央—地方行动文书链。核心取证：《清德宗实录》光绪六年相关条；《筹办夷务始末》相关年卷与中外照会、条约文本。该项锚定的是命令或决策环节，命令抵达和结果仍须另外核验。来源保留：奏折、上谕、部议、军机档或条约可证明行政动作和机构立场，不自动证明所有地区、群体或对手按其意图行动。
\eventanchor{EQ-1880-ILI-NEGOTIATIONS-AFTERMATH}{「清俄围绕伊犁归还和边界条件…}光绪六年（1880年），「清俄围绕伊犁归还和边界条件持续谈判」引发的善后、执行或地方回应构成可由不同材料追踪的后续节点。核心取证：《清德宗实录》光绪六年相关条；《筹办夷务始末》相关年卷与中外照会、条约文本；相关地区的档案、志书、契约、报刊或对方材料。该项锚定的是后续追踪问题，影响范围须按地区与群体分别判断。来源保留：战后、改革后或条约后的记忆、地方记录和官方总结常具有事后选择性；后续影响须按地点、群体和时间分层，不作整体化推断。
\eventanchor{EQ-1882-IMO-INCIDENT-PRELUDE}{「朝鲜壬午兵变后清军入朝，日…}光绪八年（1882年），「朝鲜壬午兵变后清军入朝，日本也扩大影响」之前的条件、信息和争议已通过中央与地方材料显现，构成该事件可追踪的前史节点。核心取证：《清德宗实录》光绪八年相关条；《筹办夷务始末》相关年卷与中外照会、条约文本；同时期地方档案、地方志、日记或对方材料应按具体地区补检。该项锚定的是前史条件与信息背景，不把条件直接写成唯一因果。来源保留：官修编年和地方材料的记录密度与立场并不相同；本行不将条件、传闻、呈报或后见解释直接等同为确定因果。
\eventanchor{EQ-1882-IMO-INCIDENT-DECISION}{围绕「朝鲜壬午兵变后清军入朝…}光绪八年（1882年），围绕「朝鲜壬午兵变后清军入朝，日本也扩大影响」的命令、调度、照会或呈报形成独立的中央—地方行动文书链。核心取证：《清德宗实录》光绪八年相关条；《筹办夷务始末》相关年卷与中外照会、条约文本。该项锚定的是命令或决策环节，命令抵达和结果仍须另外核验。来源保留：奏折、上谕、部议、军机档或条约可证明行政动作和机构立场，不自动证明所有地区、群体或对手按其意图行动。
\eventanchor{EQ-1882-IMO-INCIDENT-AFTERMATH}{「朝鲜壬午兵变后清军入朝，日…}光绪八年（1882年），「朝鲜壬午兵变后清军入朝，日本也扩大影响」引发的善后、执行或地方回应构成可由不同材料追踪的后续节点。核心取证：《清德宗实录》光绪八年相关条；《筹办夷务始末》相关年卷与中外照会、条约文本；相关地区的档案、志书、契约、报刊或对方材料。该项锚定的是后续追踪问题，影响范围须按地区与群体分别判断。来源保留：战后、改革后或条约后的记忆、地方记录和官方总结常具有事后选择性；后续影响须按地点、群体和时间分层，不作整体化推断。
\eventanchor{EQ-1883-SINO-FRENCH-WAR-BEGIN-PRELUDE}{「中法因越南北部和宗藩关系冲…}光绪九年（1883年），「中法因越南北部和宗藩关系冲突升级为战争」之前的条件、信息和争议已通过中央与地方材料显现，构成该事件可追踪的前史节点。核心取证：《清德宗实录》光绪九年相关条；军机处档、战事方略及地方奏报；同时期地方档案、地方志、日记或对方材料应按具体地区补检。该项锚定的是前史条件与信息背景，不把条件直接写成唯一因果。来源保留：官修编年和地方材料的记录密度与立场并不相同；本行不将条件、传闻、呈报或后见解释直接等同为确定因果。
\eventanchor{EQ-1883-SINO-FRENCH-WAR-BEGIN-DECISION}{围绕「中法因越南北部和宗藩关…}光绪九年（1883年），围绕「中法因越南北部和宗藩关系冲突升级为战争」的命令、调度、照会或呈报形成独立的中央—地方行动文书链。核心取证：《清德宗实录》光绪九年相关条；军机处档、战事方略及地方奏报。该项锚定的是命令或决策环节，命令抵达和结果仍须另外核验。来源保留：奏折、上谕、部议、军机档或条约可证明行政动作和机构立场，不自动证明所有地区、群体或对手按其意图行动。
\eventanchor{EQ-1883-SINO-FRENCH-WAR-BEGIN-AFTERMATH}{「中法因越南北部和宗藩关系冲…}光绪九年（1883年），「中法因越南北部和宗藩关系冲突升级为战争」引发的善后、执行或地方回应构成可由不同材料追踪的后续节点。核心取证：《清德宗实录》光绪九年相关条；军机处档、战事方略及地方奏报；相关地区的档案、志书、契约、报刊或对方材料。该项锚定的是后续追踪问题，影响范围须按地区与群体分别判断。来源保留：战后、改革后或条约后的记忆、地方记录和官方总结常具有事后选择性；后续影响须按地点、群体和时间分层，不作整体化推断。
\eventanchor{EQ-1884-FUZHOU-NAVAL-BATTLE-PRELUDE}{「马江海战中福建水师遭法国舰…}光绪十年（1884年），「马江海战中福建水师遭法国舰队重创」之前的条件、信息和争议已通过中央与地方材料显现，构成该事件可追踪的前史节点。核心取证：《清德宗实录》光绪十年相关条；军机处档、战事方略及地方奏报；同时期地方档案、地方志、日记或对方材料应按具体地区补检。该项锚定的是前史条件与信息背景，不把条件直接写成唯一因果。来源保留：官修编年和地方材料的记录密度与立场并不相同；本行不将条件、传闻、呈报或后见解释直接等同为确定因果。
\eventanchor{EQ-1884-FUZHOU-NAVAL-BATTLE-DECISION}{围绕「马江海战中福建水师遭法…}光绪十年（1884年），围绕「马江海战中福建水师遭法国舰队重创」的命令、调度、照会或呈报形成独立的中央—地方行动文书链。核心取证：《清德宗实录》光绪十年相关条；军机处档、战事方略及地方奏报。该项锚定的是命令或决策环节，命令抵达和结果仍须另外核验。来源保留：奏折、上谕、部议、军机档或条约可证明行政动作和机构立场，不自动证明所有地区、群体或对手按其意图行动。
\eventanchor{EQ-1884-FUZHOU-NAVAL-BATTLE-AFTERMATH}{「马江海战中福建水师遭法国舰…}光绪十年（1884年），「马江海战中福建水师遭法国舰队重创」引发的善后、执行或地方回应构成可由不同材料追踪的后续节点。核心取证：《清德宗实录》光绪十年相关条；军机处档、战事方略及地方奏报；相关地区的档案、志书、契约、报刊或对方材料。该项锚定的是后续追踪问题，影响范围须按地区与群体分别判断。来源保留：战后、改革后或条约后的记忆、地方记录和官方总结常具有事后选择性；后续影响须按地点、群体和时间分层，不作整体化推断。
\eventanchor{EQ-1885-TIANJIN-CHINA-FRANCE-TREATY-PRELUDE}{「《中法新约》确认越南脱离清…}光绪十一年（1885年），「《中法新约》确认越南脱离清朝宗藩体系，战争结束」之前的条件、信息和争议已通过中央与地方材料显现，构成该事件可追踪的前史节点。核心取证：《清德宗实录》光绪十一年相关条；《筹办夷务始末》相关年卷与中外照会、条约文本；同时期地方档案、地方志、日记或对方材料应按具体地区补检。该项锚定的是前史条件与信息背景，不把条件直接写成唯一因果。来源保留：官修编年和地方材料的记录密度与立场并不相同；本行不将条件、传闻、呈报或后见解释直接等同为确定因果。
\eventanchor{EQ-1885-TIANJIN-CHINA-FRANCE-TREATY-DECISION}{围绕「《中法新约》确认越南脱…}光绪十一年（1885年），围绕「《中法新约》确认越南脱离清朝宗藩体系，战争结束」的命令、调度、照会或呈报形成独立的中央—地方行动文书链。核心取证：《清德宗实录》光绪十一年相关条；《筹办夷务始末》相关年卷与中外照会、条约文本。该项锚定的是命令或决策环节，命令抵达和结果仍须另外核验。来源保留：奏折、上谕、部议、军机档或条约可证明行政动作和机构立场，不自动证明所有地区、群体或对手按其意图行动。
\eventanchor{EQ-1885-TIANJIN-CHINA-FRANCE-TREATY-AFTERMATH}{「《中法新约》确认越南脱离清…}光绪十一年（1885年），「《中法新约》确认越南脱离清朝宗藩体系，战争结束」引发的善后、执行或地方回应构成可由不同材料追踪的后续节点。核心取证：《清德宗实录》光绪十一年相关条；《筹办夷务始末》相关年卷与中外照会、条约文本；相关地区的档案、志书、契约、报刊或对方材料。该项锚定的是后续追踪问题，影响范围须按地区与群体分别判断。来源保留：战后、改革后或条约后的记忆、地方记录和官方总结常具有事后选择性；后续影响须按地点、群体和时间分层，不作整体化推断。
\eventanchor{EQ-1885-TAIWAN-PROVINCE-PRELUDE}{「清廷决定将台湾建省，加强行…}光绪十一年（1885年），「清廷决定将台湾建省，加强行政、军防和基础设施经营」之前的条件、信息和争议已通过中央与地方材料显现，构成该事件可追踪的前史节点。核心取证：《清德宗实录》光绪十一年相关条；宫中朱批奏折、内阁大库题本与《清史稿》相关志传；同时期地方档案、地方志、日记或对方材料应按具体地区补检。该项锚定的是前史条件与信息背景，不把条件直接写成唯一因果。来源保留：官修编年和地方材料的记录密度与立场并不相同；本行不将条件、传闻、呈报或后见解释直接等同为确定因果。
\eventanchor{EQ-1885-TAIWAN-PROVINCE-DECISION}{围绕「清廷决定将台湾建省，加…}光绪十一年（1885年），围绕「清廷决定将台湾建省，加强行政、军防和基础设施经营」的命令、调度、照会或呈报形成独立的中央—地方行动文书链。核心取证：《清德宗实录》光绪十一年相关条；宫中朱批奏折、内阁大库题本与《清史稿》相关志传。该项锚定的是命令或决策环节，命令抵达和结果仍须另外核验。来源保留：奏折、上谕、部议、军机档或条约可证明行政动作和机构立场，不自动证明所有地区、群体或对手按其意图行动。
\eventanchor{EQ-1885-TAIWAN-PROVINCE-AFTERMATH}{「清廷决定将台湾建省，加强行…}光绪十一年（1885年），「清廷决定将台湾建省，加强行政、军防和基础设施经营」引发的善后、执行或地方回应构成可由不同材料追踪的后续节点。核心取证：《清德宗实录》光绪十一年相关条；宫中朱批奏折、内阁大库题本与《清史稿》相关志传；相关地区的档案、志书、契约、报刊或对方材料。该项锚定的是后续追踪问题，影响范围须按地区与群体分别判断。来源保留：战后、改革后或条约后的记忆、地方记录和官方总结常具有事后选择性；后续影响须按地点、群体和时间分层，不作整体化推断。
\eventanchor{EQ-1886-BEIYANG-NAVY-EXPANSION-PRELUDE}{「北洋海军和沿海防务建设持续…}光绪十二年（1886年），「北洋海军和沿海防务建设持续推进」之前的条件、信息和争议已通过中央与地方材料显现，构成该事件可追踪的前史节点。核心取证：《清德宗实录》光绪十二年相关条；军机处档、战事方略及地方奏报；同时期地方档案、地方志、日记或对方材料应按具体地区补检。该项锚定的是前史条件与信息背景，不把条件直接写成唯一因果。来源保留：官修编年和地方材料的记录密度与立场并不相同；本行不将条件、传闻、呈报或后见解释直接等同为确定因果。
\eventanchor{EQ-1886-BEIYANG-NAVY-EXPANSION-DECISION}{围绕「北洋海军和沿海防务建设…}光绪十二年（1886年），围绕「北洋海军和沿海防务建设持续推进」的命令、调度、照会或呈报形成独立的中央—地方行动文书链。核心取证：《清德宗实录》光绪十二年相关条；军机处档、战事方略及地方奏报。该项锚定的是命令或决策环节，命令抵达和结果仍须另外核验。来源保留：奏折、上谕、部议、军机档或条约可证明行政动作和机构立场，不自动证明所有地区、群体或对手按其意图行动。
\eventanchor{EQ-1886-BEIYANG-NAVY-EXPANSION-AFTERMATH}{「北洋海军和沿海防务建设持续…}光绪十二年（1886年），「北洋海军和沿海防务建设持续推进」引发的善后、执行或地方回应构成可由不同材料追踪的后续节点。核心取证：《清德宗实录》光绪十二年相关条；军机处档、战事方略及地方奏报；相关地区的档案、志书、契约、报刊或对方材料。该项锚定的是后续追踪问题，影响范围须按地区与群体分别判断。来源保留：战后、改革后或条约后的记忆、地方记录和官方总结常具有事后选择性；后续影响须按地点、群体和时间分层，不作整体化推断。
\eventanchor{EQ-1888-NAVAL-AND-ARMY-REFORM-PRELUDE}{「清廷讨论海军、陆军训练和军…}光绪十四年（1888年），「清廷讨论海军、陆军训练和军工体系的经费与制度问题」之前的条件、信息和争议已通过中央与地方材料显现，构成该事件可追踪的前史节点。核心取证：《清德宗实录》光绪十四年相关条；宫中朱批奏折、内阁大库题本与《清史稿》相关志传；同时期地方档案、地方志、日记或对方材料应按具体地区补检。该项锚定的是前史条件与信息背景，不把条件直接写成唯一因果。来源保留：官修编年和地方材料的记录密度与立场并不相同；本行不将条件、传闻、呈报或后见解释直接等同为确定因果。
\eventanchor{EQ-1888-NAVAL-AND-ARMY-REFORM-DECISION}{围绕「清廷讨论海军、陆军训练…}光绪十四年（1888年），围绕「清廷讨论海军、陆军训练和军工体系的经费与制度问题」的命令、调度、照会或呈报形成独立的中央—地方行动文书链。核心取证：《清德宗实录》光绪十四年相关条；宫中朱批奏折、内阁大库题本与《清史稿》相关志传。该项锚定的是命令或决策环节，命令抵达和结果仍须另外核验。来源保留：奏折、上谕、部议、军机档或条约可证明行政动作和机构立场，不自动证明所有地区、群体或对手按其意图行动。
\eventanchor{EQ-1888-NAVAL-AND-ARMY-REFORM-AFTERMATH}{「清廷讨论海军、陆军训练和军…}光绪十四年（1888年），「清廷讨论海军、陆军训练和军工体系的经费与制度问题」引发的善后、执行或地方回应构成可由不同材料追踪的后续节点。核心取证：《清德宗实录》光绪十四年相关条；宫中朱批奏折、内阁大库题本与《清史稿》相关志传；相关地区的档案、志书、契约、报刊或对方材料。该项锚定的是后续追踪问题，影响范围须按地区与群体分别判断。来源保留：战后、改革后或条约后的记忆、地方记录和官方总结常具有事后选择性；后续影响须按地点、群体和时间分层，不作整体化推断。
\eventanchor{EQ-1889-GUANGXU-MARRIAGE-PRELUDE}{「光绪帝大婚并开始形式上的亲…}光绪十五年（1889年），「光绪帝大婚并开始形式上的亲政安排」之前的条件、信息和争议已通过中央与地方材料显现，构成该事件可追踪的前史节点。核心取证：《清德宗实录》光绪十五年相关条；宫中朱批奏折、内阁大库题本与《清史稿》相关志传；同时期地方档案、地方志、日记或对方材料应按具体地区补检。该项锚定的是前史条件与信息背景，不把条件直接写成唯一因果。来源保留：官修编年和地方材料的记录密度与立场并不相同；本行不将条件、传闻、呈报或后见解释直接等同为确定因果。
\eventanchor{EQ-1889-GUANGXU-MARRIAGE-DECISION}{围绕「光绪帝大婚并开始形式上…}光绪十五年（1889年），围绕「光绪帝大婚并开始形式上的亲政安排」的命令、调度、照会或呈报形成独立的中央—地方行动文书链。核心取证：《清德宗实录》光绪十五年相关条；宫中朱批奏折、内阁大库题本与《清史稿》相关志传。该项锚定的是命令或决策环节，命令抵达和结果仍须另外核验。来源保留：奏折、上谕、部议、军机档或条约可证明行政动作和机构立场，不自动证明所有地区、群体或对手按其意图行动。
\eventanchor{EQ-1889-GUANGXU-MARRIAGE-AFTERMATH}{「光绪帝大婚并开始形式上的亲…}光绪十五年（1889年），「光绪帝大婚并开始形式上的亲政安排」引发的善后、执行或地方回应构成可由不同材料追踪的后续节点。核心取证：《清德宗实录》光绪十五年相关条；宫中朱批奏折、内阁大库题本与《清史稿》相关志传；相关地区的档案、志书、契约、报刊或对方材料。该项锚定的是后续追踪问题，影响范围须按地区与群体分别判断。来源保留：战后、改革后或条约后的记忆、地方记录和官方总结常具有事后选择性；后续影响须按地点、群体和时间分层，不作整体化推断。
\subsection{1890—1899}
\eventanchor{EQ-1890-KOREAN-POLITICAL-CONFLICT-PRELUDE}{「朝鲜内政纷争持续牵动清日两…}光绪十六年（1890年），「朝鲜内政纷争持续牵动清日两国在半岛的影响力竞争」之前的条件、信息和争议已通过中央与地方材料显现，构成该事件可追踪的前史节点。核心取证：《清德宗实录》光绪十六年相关条；《筹办夷务始末》相关年卷与中外照会、条约文本；同时期地方档案、地方志、日记或对方材料应按具体地区补检。该项锚定的是前史条件与信息背景，不把条件直接写成唯一因果。来源保留：官修编年和地方材料的记录密度与立场并不相同；本行不将条件、传闻、呈报或后见解释直接等同为确定因果。
\eventanchor{EQ-1890-KOREAN-POLITICAL-CONFLICT-DECISION}{围绕「朝鲜内政纷争持续牵动清…}光绪十六年（1890年），围绕「朝鲜内政纷争持续牵动清日两国在半岛的影响力竞争」的命令、调度、照会或呈报形成独立的中央—地方行动文书链。核心取证：《清德宗实录》光绪十六年相关条；《筹办夷务始末》相关年卷与中外照会、条约文本。该项锚定的是命令或决策环节，命令抵达和结果仍须另外核验。来源保留：奏折、上谕、部议、军机档或条约可证明行政动作和机构立场，不自动证明所有地区、群体或对手按其意图行动。
\eventanchor{EQ-1890-KOREAN-POLITICAL-CONFLICT-AFTERMATH}{「朝鲜内政纷争持续牵动清日两…}光绪十六年（1890年），「朝鲜内政纷争持续牵动清日两国在半岛的影响力竞争」引发的善后、执行或地方回应构成可由不同材料追踪的后续节点。核心取证：《清德宗实录》光绪十六年相关条；《筹办夷务始末》相关年卷与中外照会、条约文本；相关地区的档案、志书、契约、报刊或对方材料。该项锚定的是后续追踪问题，影响范围须按地区与群体分别判断。来源保留：战后、改革后或条约后的记忆、地方记录和官方总结常具有事后选择性；后续影响须按地点、群体和时间分层，不作整体化推断。
\eventanchor{EQ-1892-RAILWAY-DEBATE-PRELUDE}{「官员、商人和士绅围绕铁路修…}光绪十八年（1892年），「官员、商人和士绅围绕铁路修筑、主权与资金展开争论」之前的条件、信息和争议已通过中央与地方材料显现，构成该事件可追踪的前史节点。核心取证：《清德宗实录》光绪十八年相关条；户部奏销、盐政、河工或海关档案；同时期地方档案、地方志、日记或对方材料应按具体地区补检。该项锚定的是前史条件与信息背景，不把条件直接写成唯一因果。来源保留：官修编年和地方材料的记录密度与立场并不相同；本行不将条件、传闻、呈报或后见解释直接等同为确定因果。
\eventanchor{EQ-1892-RAILWAY-DEBATE-DECISION}{围绕「官员、商人和士绅围绕铁…}光绪十八年（1892年），围绕「官员、商人和士绅围绕铁路修筑、主权与资金展开争论」的命令、调度、照会或呈报形成独立的中央—地方行动文书链。核心取证：《清德宗实录》光绪十八年相关条；户部奏销、盐政、河工或海关档案。该项锚定的是命令或决策环节，命令抵达和结果仍须另外核验。来源保留：奏折、上谕、部议、军机档或条约可证明行政动作和机构立场，不自动证明所有地区、群体或对手按其意图行动。
\eventanchor{EQ-1892-RAILWAY-DEBATE-AFTERMATH}{「官员、商人和士绅围绕铁路修…}光绪十八年（1892年），「官员、商人和士绅围绕铁路修筑、主权与资金展开争论」引发的善后、执行或地方回应构成可由不同材料追踪的后续节点。核心取证：《清德宗实录》光绪十八年相关条；户部奏销、盐政、河工或海关档案；相关地区的档案、志书、契约、报刊或对方材料。该项锚定的是后续追踪问题，影响范围须按地区与群体分别判断。来源保留：战后、改革后或条约后的记忆、地方记录和官方总结常具有事后选择性；后续影响须按地点、群体和时间分层，不作整体化推断。
\eventanchor{EQ-1893-KOREAN-REFORM-CRISIS-PRELUDE}{「朝鲜改革与政治危机使清日关…}光绪十九年（1893年），「朝鲜改革与政治危机使清日关系持续紧张」之前的条件、信息和争议已通过中央与地方材料显现，构成该事件可追踪的前史节点。核心取证：《清德宗实录》光绪十九年相关条；《筹办夷务始末》相关年卷与中外照会、条约文本；同时期地方档案、地方志、日记或对方材料应按具体地区补检。该项锚定的是前史条件与信息背景，不把条件直接写成唯一因果。来源保留：官修编年和地方材料的记录密度与立场并不相同；本行不将条件、传闻、呈报或后见解释直接等同为确定因果。
\eventanchor{EQ-1893-KOREAN-REFORM-CRISIS-DECISION}{围绕「朝鲜改革与政治危机使清…}光绪十九年（1893年），围绕「朝鲜改革与政治危机使清日关系持续紧张」的命令、调度、照会或呈报形成独立的中央—地方行动文书链。核心取证：《清德宗实录》光绪十九年相关条；《筹办夷务始末》相关年卷与中外照会、条约文本。该项锚定的是命令或决策环节，命令抵达和结果仍须另外核验。来源保留：奏折、上谕、部议、军机档或条约可证明行政动作和机构立场，不自动证明所有地区、群体或对手按其意图行动。
\eventanchor{EQ-1893-KOREAN-REFORM-CRISIS-AFTERMATH}{「朝鲜改革与政治危机使清日关…}光绪十九年（1893年），「朝鲜改革与政治危机使清日关系持续紧张」引发的善后、执行或地方回应构成可由不同材料追踪的后续节点。核心取证：《清德宗实录》光绪十九年相关条；《筹办夷务始末》相关年卷与中外照会、条约文本；相关地区的档案、志书、契约、报刊或对方材料。该项锚定的是后续追踪问题，影响范围须按地区与群体分别判断。来源保留：战后、改革后或条约后的记忆、地方记录和官方总结常具有事后选择性；后续影响须按地点、群体和时间分层，不作整体化推断。
\eventanchor{EQ-1894-DONGHAK-PEASANT-WAR-PRELUDE}{「东学农民战争促使清日两国派…}光绪二十年（1894年），「东学农民战争促使清日两国派兵入朝，危机迅速国际化」之前的条件、信息和争议已通过中央与地方材料显现，构成该事件可追踪的前史节点。核心取证：《清德宗实录》光绪二十年相关条；军机处档、战事方略及地方奏报；同时期地方档案、地方志、日记或对方材料应按具体地区补检。该项锚定的是前史条件与信息背景，不把条件直接写成唯一因果。来源保留：官修编年和地方材料的记录密度与立场并不相同；本行不将条件、传闻、呈报或后见解释直接等同为确定因果。
\eventanchor{EQ-1894-DONGHAK-PEASANT-WAR-DECISION}{围绕「东学农民战争促使清日两…}光绪二十年（1894年），围绕「东学农民战争促使清日两国派兵入朝，危机迅速国际化」的命令、调度、照会或呈报形成独立的中央—地方行动文书链。核心取证：《清德宗实录》光绪二十年相关条；军机处档、战事方略及地方奏报。该项锚定的是命令或决策环节，命令抵达和结果仍须另外核验。来源保留：奏折、上谕、部议、军机档或条约可证明行政动作和机构立场，不自动证明所有地区、群体或对手按其意图行动。
\eventanchor{EQ-1894-DONGHAK-PEASANT-WAR-AFTERMATH}{「东学农民战争促使清日两国派…}光绪二十年（1894年），「东学农民战争促使清日两国派兵入朝，危机迅速国际化」引发的善后、执行或地方回应构成可由不同材料追踪的后续节点。核心取证：《清德宗实录》光绪二十年相关条；军机处档、战事方略及地方奏报；相关地区的档案、志书、契约、报刊或对方材料。该项锚定的是后续追踪问题，影响范围须按地区与群体分别判断。来源保留：战后、改革后或条约后的记忆、地方记录和官方总结常具有事后选择性；后续影响须按地点、群体和时间分层，不作整体化推断。
\eventanchor{EQ-1894-YALU-NAVAL-BATTLE-PRELUDE}{「黄海海战后北洋舰队受创，海…}光绪二十年（1894年），「黄海海战后北洋舰队受创，海军力量对比进一步恶化」之前的条件、信息和争议已通过中央与地方材料显现，构成该事件可追踪的前史节点。核心取证：《清德宗实录》光绪二十年相关条；军机处档、战事方略及地方奏报；同时期地方档案、地方志、日记或对方材料应按具体地区补检。该项锚定的是前史条件与信息背景，不把条件直接写成唯一因果。来源保留：官修编年和地方材料的记录密度与立场并不相同；本行不将条件、传闻、呈报或后见解释直接等同为确定因果。
\eventanchor{EQ-1894-YALU-NAVAL-BATTLE-DECISION}{围绕「黄海海战后北洋舰队受创…}光绪二十年（1894年），围绕「黄海海战后北洋舰队受创，海军力量对比进一步恶化」的命令、调度、照会或呈报形成独立的中央—地方行动文书链。核心取证：《清德宗实录》光绪二十年相关条；军机处档、战事方略及地方奏报。该项锚定的是命令或决策环节，命令抵达和结果仍须另外核验。来源保留：奏折、上谕、部议、军机档或条约可证明行政动作和机构立场，不自动证明所有地区、群体或对手按其意图行动。
\eventanchor{EQ-1894-YALU-NAVAL-BATTLE-AFTERMATH}{「黄海海战后北洋舰队受创，海…}光绪二十年（1894年），「黄海海战后北洋舰队受创，海军力量对比进一步恶化」引发的善后、执行或地方回应构成可由不同材料追踪的后续节点。核心取证：《清德宗实录》光绪二十年相关条；军机处档、战事方略及地方奏报；相关地区的档案、志书、契约、报刊或对方材料。该项锚定的是后续追踪问题，影响范围须按地区与群体分别判断。来源保留：战后、改革后或条约后的记忆、地方记录和官方总结常具有事后选择性；后续影响须按地点、群体和时间分层，不作整体化推断。
\eventanchor{EQ-1895-SHIMONOSEKI-TREATY-PRELUDE}{「《马关条约》签订，清割让台…}光绪二十一年（1895年），「《马关条约》签订，清割让台湾、澎湖并承认朝鲜独立，赔款与通商条件加重」之前的条件、信息和争议已通过中央与地方材料显现，构成该事件可追踪的前史节点。核心取证：《清德宗实录》光绪二十一年相关条；《筹办夷务始末》相关年卷与中外照会、条约文本；同时期地方档案、地方志、日记或对方材料应按具体地区补检。该项锚定的是前史条件与信息背景，不把条件直接写成唯一因果。来源保留：官修编年和地方材料的记录密度与立场并不相同；本行不将条件、传闻、呈报或后见解释直接等同为确定因果。
\eventanchor{EQ-1895-SHIMONOSEKI-TREATY-DECISION}{围绕「《马关条约》签订，清割…}光绪二十一年（1895年），围绕「《马关条约》签订，清割让台湾、澎湖并承认朝鲜独立，赔款与通商条件加重」的命令、调度、照会或呈报形成独立的中央—地方行动文书链。核心取证：《清德宗实录》光绪二十一年相关条；《筹办夷务始末》相关年卷与中外照会、条约文本。该项锚定的是命令或决策环节，命令抵达和结果仍须另外核验。来源保留：奏折、上谕、部议、军机档或条约可证明行政动作和机构立场，不自动证明所有地区、群体或对手按其意图行动。
\eventanchor{EQ-1895-SHIMONOSEKI-TREATY-AFTERMATH}{「《马关条约》签订，清割让台…}光绪二十一年（1895年），「《马关条约》签订，清割让台湾、澎湖并承认朝鲜独立，赔款与通商条件加重」引发的善后、执行或地方回应构成可由不同材料追踪的后续节点。核心取证：《清德宗实录》光绪二十一年相关条；《筹办夷务始末》相关年卷与中外照会、条约文本；相关地区的档案、志书、契约、报刊或对方材料。该项锚定的是后续追踪问题，影响范围须按地区与群体分别判断。来源保留：战后、改革后或条约后的记忆、地方记录和官方总结常具有事后选择性；后续影响须按地点、群体和时间分层，不作整体化推断。
\eventanchor{EQ-1895-TAIWAN-REPUBLIC-RESISTANCE-PRELUDE}{「台湾割让后地方力量建立台湾…}光绪二十一年（1895年），「台湾割让后地方力量建立台湾民主国并抵抗日军登陆」之前的条件、信息和争议已通过中央与地方材料显现，构成该事件可追踪的前史节点。核心取证：《清德宗实录》光绪二十一年相关条；军机处档、战事方略及地方奏报；同时期地方档案、地方志、日记或对方材料应按具体地区补检。该项锚定的是前史条件与信息背景，不把条件直接写成唯一因果。来源保留：官修编年和地方材料的记录密度与立场并不相同；本行不将条件、传闻、呈报或后见解释直接等同为确定因果。
\eventanchor{EQ-1895-TAIWAN-REPUBLIC-RESISTANCE-DECISION}{围绕「台湾割让后地方力量建立…}光绪二十一年（1895年），围绕「台湾割让后地方力量建立台湾民主国并抵抗日军登陆」的命令、调度、照会或呈报形成独立的中央—地方行动文书链。核心取证：《清德宗实录》光绪二十一年相关条；军机处档、战事方略及地方奏报。该项锚定的是命令或决策环节，命令抵达和结果仍须另外核验。来源保留：奏折、上谕、部议、军机档或条约可证明行政动作和机构立场，不自动证明所有地区、群体或对手按其意图行动。
\eventanchor{EQ-1895-TAIWAN-REPUBLIC-RESISTANCE-AFTERMATH}{「台湾割让后地方力量建立台湾…}光绪二十一年（1895年），「台湾割让后地方力量建立台湾民主国并抵抗日军登陆」引发的善后、执行或地方回应构成可由不同材料追踪的后续节点。核心取证：《清德宗实录》光绪二十一年相关条；军机处档、战事方略及地方奏报；相关地区的档案、志书、契约、报刊或对方材料。该项锚定的是后续追踪问题，影响范围须按地区与群体分别判断。来源保留：战后、改革后或条约后的记忆、地方记录和官方总结常具有事后选择性；后续影响须按地点、群体和时间分层，不作整体化推断。
\eventanchor{EQ-1896-LI-LOBANOV-TREATY-PRELUDE}{「《中俄密约》及东清铁路安排…}光绪二十二年（1896年），「《中俄密约》及东清铁路安排加强俄国在东北的影响」之前的条件、信息和争议已通过中央与地方材料显现，构成该事件可追踪的前史节点。核心取证：《清德宗实录》光绪二十二年相关条；《筹办夷务始末》相关年卷与中外照会、条约文本；同时期地方档案、地方志、日记或对方材料应按具体地区补检。该项锚定的是前史条件与信息背景，不把条件直接写成唯一因果。来源保留：官修编年和地方材料的记录密度与立场并不相同；本行不将条件、传闻、呈报或后见解释直接等同为确定因果。
\eventanchor{EQ-1896-LI-LOBANOV-TREATY-DECISION}{围绕「《中俄密约》及东清铁路…}光绪二十二年（1896年），围绕「《中俄密约》及东清铁路安排加强俄国在东北的影响」的命令、调度、照会或呈报形成独立的中央—地方行动文书链。核心取证：《清德宗实录》光绪二十二年相关条；《筹办夷务始末》相关年卷与中外照会、条约文本。该项锚定的是命令或决策环节，命令抵达和结果仍须另外核验。来源保留：奏折、上谕、部议、军机档或条约可证明行政动作和机构立场，不自动证明所有地区、群体或对手按其意图行动。
\eventanchor{EQ-1896-LI-LOBANOV-TREATY-AFTERMATH}{「《中俄密约》及东清铁路安排…}光绪二十二年（1896年），「《中俄密约》及东清铁路安排加强俄国在东北的影响」引发的善后、执行或地方回应构成可由不同材料追踪的后续节点。核心取证：《清德宗实录》光绪二十二年相关条；《筹办夷务始末》相关年卷与中外照会、条约文本；相关地区的档案、志书、契约、报刊或对方材料。该项锚定的是后续追踪问题，影响范围须按地区与群体分别判断。来源保留：战后、改革后或条约后的记忆、地方记录和官方总结常具有事后选择性；后续影响须按地点、群体和时间分层，不作整体化推断。
\eventanchor{EQ-1896-MODERN-EDUCATION-EXPANSION-PRELUDE}{「甲午战后新式学堂、译书与留…}光绪二十二年（1896年），「甲午战后新式学堂、译书与留学讨论加速」之前的条件、信息和争议已通过中央与地方材料显现，构成该事件可追踪的前史节点。核心取证：《清德宗实录》光绪二十二年相关条；内阁档、文集或报刊相关记录；同时期地方档案、地方志、日记或对方材料应按具体地区补检。该项锚定的是前史条件与信息背景，不把条件直接写成唯一因果。来源保留：官修编年和地方材料的记录密度与立场并不相同；本行不将条件、传闻、呈报或后见解释直接等同为确定因果。
\eventanchor{EQ-1896-MODERN-EDUCATION-EXPANSION-DECISION}{围绕「甲午战后新式学堂、译书…}光绪二十二年（1896年），围绕「甲午战后新式学堂、译书与留学讨论加速」的命令、调度、照会或呈报形成独立的中央—地方行动文书链。核心取证：《清德宗实录》光绪二十二年相关条；内阁档、文集或报刊相关记录。该项锚定的是命令或决策环节，命令抵达和结果仍须另外核验。来源保留：奏折、上谕、部议、军机档或条约可证明行政动作和机构立场，不自动证明所有地区、群体或对手按其意图行动。
\eventanchor{EQ-1896-MODERN-EDUCATION-EXPANSION-AFTERMATH}{「甲午战后新式学堂、译书与留…}光绪二十二年（1896年），「甲午战后新式学堂、译书与留学讨论加速」引发的善后、执行或地方回应构成可由不同材料追踪的后续节点。核心取证：《清德宗实录》光绪二十二年相关条；内阁档、文集或报刊相关记录；相关地区的档案、志书、契约、报刊或对方材料。该项锚定的是后续追踪问题，影响范围须按地区与群体分别判断。来源保留：战后、改革后或条约后的记忆、地方记录和官方总结常具有事后选择性；后续影响须按地点、群体和时间分层，不作整体化推断。
\eventanchor{EQ-1897-JIAOZHOU-BAY-OCCUPATION-PRELUDE}{「德国以巨野教案为由占领胶州…}光绪二十三年（1897年），「德国以巨野教案为由占领胶州湾，清廷面临新的租借压力」之前的条件、信息和争议已通过中央与地方材料显现，构成该事件可追踪的前史节点。核心取证：《清德宗实录》光绪二十三年相关条；《筹办夷务始末》相关年卷与中外照会、条约文本；同时期地方档案、地方志、日记或对方材料应按具体地区补检。该项锚定的是前史条件与信息背景，不把条件直接写成唯一因果。来源保留：官修编年和地方材料的记录密度与立场并不相同；本行不将条件、传闻、呈报或后见解释直接等同为确定因果。
\eventanchor{EQ-1897-JIAOZHOU-BAY-OCCUPATION-DECISION}{围绕「德国以巨野教案为由占领…}光绪二十三年（1897年），围绕「德国以巨野教案为由占领胶州湾，清廷面临新的租借压力」的命令、调度、照会或呈报形成独立的中央—地方行动文书链。核心取证：《清德宗实录》光绪二十三年相关条；《筹办夷务始末》相关年卷与中外照会、条约文本。该项锚定的是命令或决策环节，命令抵达和结果仍须另外核验。来源保留：奏折、上谕、部议、军机档或条约可证明行政动作和机构立场，不自动证明所有地区、群体或对手按其意图行动。
\eventanchor{EQ-1897-JIAOZHOU-BAY-OCCUPATION-AFTERMATH}{「德国以巨野教案为由占领胶州…}光绪二十三年（1897年），「德国以巨野教案为由占领胶州湾，清廷面临新的租借压力」引发的善后、执行或地方回应构成可由不同材料追踪的后续节点。核心取证：《清德宗实录》光绪二十三年相关条；《筹办夷务始末》相关年卷与中外照会、条约文本；相关地区的档案、志书、契约、报刊或对方材料。该项锚定的是后续追踪问题，影响范围须按地区与群体分别判断。来源保留：战后、改革后或条约后的记忆、地方记录和官方总结常具有事后选择性；后续影响须按地点、群体和时间分层，不作整体化推断。
\eventanchor{EQ-1898-LEASED-TERRITORIES-PRELUDE}{「俄、英、法、德相继取得租借…}光绪二十四年（1898年），「俄、英、法、德相继取得租借地或势力范围安排」之前的条件、信息和争议已通过中央与地方材料显现，构成该事件可追踪的前史节点。核心取证：《清德宗实录》光绪二十四年相关条；《筹办夷务始末》相关年卷与中外照会、条约文本；同时期地方档案、地方志、日记或对方材料应按具体地区补检。该项锚定的是前史条件与信息背景，不把条件直接写成唯一因果。来源保留：官修编年和地方材料的记录密度与立场并不相同；本行不将条件、传闻、呈报或后见解释直接等同为确定因果。
\eventanchor{EQ-1898-LEASED-TERRITORIES-DECISION}{围绕「俄、英、法、德相继取得…}光绪二十四年（1898年），围绕「俄、英、法、德相继取得租借地或势力范围安排」的命令、调度、照会或呈报形成独立的中央—地方行动文书链。核心取证：《清德宗实录》光绪二十四年相关条；《筹办夷务始末》相关年卷与中外照会、条约文本。该项锚定的是命令或决策环节，命令抵达和结果仍须另外核验。来源保留：奏折、上谕、部议、军机档或条约可证明行政动作和机构立场，不自动证明所有地区、群体或对手按其意图行动。
\eventanchor{EQ-1898-LEASED-TERRITORIES-AFTERMATH}{「俄、英、法、德相继取得租借…}光绪二十四年（1898年），「俄、英、法、德相继取得租借地或势力范围安排」引发的善后、执行或地方回应构成可由不同材料追踪的后续节点。核心取证：《清德宗实录》光绪二十四年相关条；《筹办夷务始末》相关年卷与中外照会、条约文本；相关地区的档案、志书、契约、报刊或对方材料。该项锚定的是后续追踪问题，影响范围须按地区与群体分别判断。来源保留：战后、改革后或条约后的记忆、地方记录和官方总结常具有事后选择性；后续影响须按地点、群体和时间分层，不作整体化推断。
\eventanchor{EQ-1898-WUXU-COUP-PRELUDE}{「慈禧发动政变，光绪帝被幽禁…}光绪二十四年（1898年），「慈禧发动政变，光绪帝被幽禁，戊戌变法主要措施被撤销或改写」之前的条件、信息和争议已通过中央与地方材料显现，构成该事件可追踪的前史节点。核心取证：《清德宗实录》光绪二十四年相关条；宫中朱批奏折、内阁大库题本与《清史稿》相关志传；同时期地方档案、地方志、日记或对方材料应按具体地区补检。该项锚定的是前史条件与信息背景，不把条件直接写成唯一因果。来源保留：官修编年和地方材料的记录密度与立场并不相同；本行不将条件、传闻、呈报或后见解释直接等同为确定因果。
\eventanchor{EQ-1898-WUXU-COUP-DECISION}{围绕「慈禧发动政变，光绪帝被…}光绪二十四年（1898年），围绕「慈禧发动政变，光绪帝被幽禁，戊戌变法主要措施被撤销或改写」的命令、调度、照会或呈报形成独立的中央—地方行动文书链。核心取证：《清德宗实录》光绪二十四年相关条；宫中朱批奏折、内阁大库题本与《清史稿》相关志传。该项锚定的是命令或决策环节，命令抵达和结果仍须另外核验。来源保留：奏折、上谕、部议、军机档或条约可证明行政动作和机构立场，不自动证明所有地区、群体或对手按其意图行动。
\eventanchor{EQ-1898-WUXU-COUP-AFTERMATH}{「慈禧发动政变，光绪帝被幽禁…}光绪二十四年（1898年），「慈禧发动政变，光绪帝被幽禁，戊戌变法主要措施被撤销或改写」引发的善后、执行或地方回应构成可由不同材料追踪的后续节点。核心取证：《清德宗实录》光绪二十四年相关条；宫中朱批奏折、内阁大库题本与《清史稿》相关志传；相关地区的档案、志书、契约、报刊或对方材料。该项锚定的是后续追踪问题，影响范围须按地区与群体分别判断。来源保留：战后、改革后或条约后的记忆、地方记录和官方总结常具有事后选择性；后续影响须按地点、群体和时间分层，不作整体化推断。
\eventanchor{EQ-1899-BOXER-UPRISING-PRELUDE}{「义和团在华北扩展，反教、反…}光绪二十五年（1899年），「义和团在华北扩展，反教、反洋冲突加剧」之前的条件、信息和争议已通过中央与地方材料显现，构成该事件可追踪的前史节点。核心取证：《清德宗实录》光绪二十五年相关条；地方志、督抚奏折及地方档案；同时期地方档案、地方志、日记或对方材料应按具体地区补检。该项锚定的是前史条件与信息背景，不把条件直接写成唯一因果。来源保留：官修编年和地方材料的记录密度与立场并不相同；本行不将条件、传闻、呈报或后见解释直接等同为确定因果。
\eventanchor{EQ-1899-BOXER-UPRISING-DECISION}{围绕「义和团在华北扩展，反教…}光绪二十五年（1899年），围绕「义和团在华北扩展，反教、反洋冲突加剧」的命令、调度、照会或呈报形成独立的中央—地方行动文书链。核心取证：《清德宗实录》光绪二十五年相关条；地方志、督抚奏折及地方档案。该项锚定的是命令或决策环节，命令抵达和结果仍须另外核验。来源保留：奏折、上谕、部议、军机档或条约可证明行政动作和机构立场，不自动证明所有地区、群体或对手按其意图行动。
\eventanchor{EQ-1899-BOXER-UPRISING-AFTERMATH}{「义和团在华北扩展，反教、反…}光绪二十五年（1899年），「义和团在华北扩展，反教、反洋冲突加剧」引发的善后、执行或地方回应构成可由不同材料追踪的后续节点。核心取证：《清德宗实录》光绪二十五年相关条；地方志、督抚奏折及地方档案；相关地区的档案、志书、契约、报刊或对方材料。该项锚定的是后续追踪问题，影响范围须按地区与群体分别判断。来源保留：战后、改革后或条约后的记忆、地方记录和官方总结常具有事后选择性；后续影响须按地点、群体和时间分层，不作整体化推断。
\subsection{1900—1909}
\eventanchor{EQ-1900-BOXER-WAR-PRELUDE}{「慈禧对列强宣战，八国联军入…}光绪二十六年（1900年），「慈禧对列强宣战，八国联军入侵并占领北京」之前的条件、信息和争议已通过中央与地方材料显现，构成该事件可追踪的前史节点。核心取证：《清德宗实录》光绪二十六年相关条；军机处档、战事方略及地方奏报；同时期地方档案、地方志、日记或对方材料应按具体地区补检。该项锚定的是前史条件与信息背景，不把条件直接写成唯一因果。来源保留：官修编年和地方材料的记录密度与立场并不相同；本行不将条件、传闻、呈报或后见解释直接等同为确定因果。
\eventanchor{EQ-1900-BOXER-WAR-DECISION}{围绕「慈禧对列强宣战，八国联…}光绪二十六年（1900年），围绕「慈禧对列强宣战，八国联军入侵并占领北京」的命令、调度、照会或呈报形成独立的中央—地方行动文书链。核心取证：《清德宗实录》光绪二十六年相关条；军机处档、战事方略及地方奏报。该项锚定的是命令或决策环节，命令抵达和结果仍须另外核验。来源保留：奏折、上谕、部议、军机档或条约可证明行政动作和机构立场，不自动证明所有地区、群体或对手按其意图行动。
\eventanchor{EQ-1900-BOXER-WAR-AFTERMATH}{「慈禧对列强宣战，八国联军入…}光绪二十六年（1900年），「慈禧对列强宣战，八国联军入侵并占领北京」引发的善后、执行或地方回应构成可由不同材料追踪的后续节点。核心取证：《清德宗实录》光绪二十六年相关条；军机处档、战事方略及地方奏报；相关地区的档案、志书、契约、报刊或对方材料。该项锚定的是后续追踪问题，影响范围须按地区与群体分别判断。来源保留：战后、改革后或条约后的记忆、地方记录和官方总结常具有事后选择性；后续影响须按地点、群体和时间分层，不作整体化推断。
\eventanchor{EQ-1901-BOXER-PROTOCOL-PRELUDE}{「《辛丑条约》签订，赔款、驻…}光绪二十七年（1901年），「《辛丑条约》签订，赔款、驻军和外交条件进一步限制清廷」之前的条件、信息和争议已通过中央与地方材料显现，构成该事件可追踪的前史节点。核心取证：《清德宗实录》光绪二十七年相关条；《筹办夷务始末》相关年卷与中外照会、条约文本；同时期地方档案、地方志、日记或对方材料应按具体地区补检。该项锚定的是前史条件与信息背景，不把条件直接写成唯一因果。来源保留：官修编年和地方材料的记录密度与立场并不相同；本行不将条件、传闻、呈报或后见解释直接等同为确定因果。
\eventanchor{EQ-1901-BOXER-PROTOCOL-DECISION}{围绕「《辛丑条约》签订，赔款…}光绪二十七年（1901年），围绕「《辛丑条约》签订，赔款、驻军和外交条件进一步限制清廷」的命令、调度、照会或呈报形成独立的中央—地方行动文书链。核心取证：《清德宗实录》光绪二十七年相关条；《筹办夷务始末》相关年卷与中外照会、条约文本。该项锚定的是命令或决策环节，命令抵达和结果仍须另外核验。来源保留：奏折、上谕、部议、军机档或条约可证明行政动作和机构立场，不自动证明所有地区、群体或对手按其意图行动。
\eventanchor{EQ-1901-BOXER-PROTOCOL-AFTERMATH}{「《辛丑条约》签订，赔款、驻…}光绪二十七年（1901年），「《辛丑条约》签订，赔款、驻军和外交条件进一步限制清廷」引发的善后、执行或地方回应构成可由不同材料追踪的后续节点。核心取证：《清德宗实录》光绪二十七年相关条；《筹办夷务始末》相关年卷与中外照会、条约文本；相关地区的档案、志书、契约、报刊或对方材料。该项锚定的是后续追踪问题，影响范围须按地区与群体分别判断。来源保留：战后、改革后或条约后的记忆、地方记录和官方总结常具有事后选择性；后续影响须按地点、群体和时间分层，不作整体化推断。
\eventanchor{EQ-1901-ZONGLI-YAMEN-REPLACED-PRELUDE}{「总理衙门改为外务部，晚清外…}光绪二十七年（1901年），「总理衙门改为外务部，晚清外交机构进一步制度化」之前的条件、信息和争议已通过中央与地方材料显现，构成该事件可追踪的前史节点。核心取证：《清德宗实录》光绪二十七年相关条；宫中朱批奏折、内阁大库题本与《清史稿》相关志传；同时期地方档案、地方志、日记或对方材料应按具体地区补检。该项锚定的是前史条件与信息背景，不把条件直接写成唯一因果。来源保留：官修编年和地方材料的记录密度与立场并不相同；本行不将条件、传闻、呈报或后见解释直接等同为确定因果。
\eventanchor{EQ-1901-ZONGLI-YAMEN-REPLACED-DECISION}{围绕「总理衙门改为外务部，晚…}光绪二十七年（1901年），围绕「总理衙门改为外务部，晚清外交机构进一步制度化」的命令、调度、照会或呈报形成独立的中央—地方行动文书链。核心取证：《清德宗实录》光绪二十七年相关条；宫中朱批奏折、内阁大库题本与《清史稿》相关志传。该项锚定的是命令或决策环节，命令抵达和结果仍须另外核验。来源保留：奏折、上谕、部议、军机档或条约可证明行政动作和机构立场，不自动证明所有地区、群体或对手按其意图行动。
\eventanchor{EQ-1901-ZONGLI-YAMEN-REPLACED-AFTERMATH}{「总理衙门改为外务部，晚清外…}光绪二十七年（1901年），「总理衙门改为外务部，晚清外交机构进一步制度化」引发的善后、执行或地方回应构成可由不同材料追踪的后续节点。核心取证：《清德宗实录》光绪二十七年相关条；宫中朱批奏折、内阁大库题本与《清史稿》相关志传；相关地区的档案、志书、契约、报刊或对方材料。该项锚定的是后续追踪问题，影响范围须按地区与群体分别判断。来源保留：战后、改革后或条约后的记忆、地方记录和官方总结常具有事后选择性；后续影响须按地点、群体和时间分层，不作整体化推断。
\eventanchor{EQ-1902-CIXI-RETURN-BEIJING-PRELUDE}{「慈禧与光绪帝返回北京，战后…}光绪二十八年（1902年），「慈禧与光绪帝返回北京，战后朝廷开始推出新政措施」之前的条件、信息和争议已通过中央与地方材料显现，构成该事件可追踪的前史节点。核心取证：《清德宗实录》光绪二十八年相关条；宫中朱批奏折、内阁大库题本与《清史稿》相关志传；同时期地方档案、地方志、日记或对方材料应按具体地区补检。该项锚定的是前史条件与信息背景，不把条件直接写成唯一因果。来源保留：官修编年和地方材料的记录密度与立场并不相同；本行不将条件、传闻、呈报或后见解释直接等同为确定因果。
\eventanchor{EQ-1902-CIXI-RETURN-BEIJING-DECISION}{围绕「慈禧与光绪帝返回北京，…}光绪二十八年（1902年），围绕「慈禧与光绪帝返回北京，战后朝廷开始推出新政措施」的命令、调度、照会或呈报形成独立的中央—地方行动文书链。核心取证：《清德宗实录》光绪二十八年相关条；宫中朱批奏折、内阁大库题本与《清史稿》相关志传。该项锚定的是命令或决策环节，命令抵达和结果仍须另外核验。来源保留：奏折、上谕、部议、军机档或条约可证明行政动作和机构立场，不自动证明所有地区、群体或对手按其意图行动。
\eventanchor{EQ-1902-CIXI-RETURN-BEIJING-AFTERMATH}{「慈禧与光绪帝返回北京，战后…}光绪二十八年（1902年），「慈禧与光绪帝返回北京，战后朝廷开始推出新政措施」引发的善后、执行或地方回应构成可由不同材料追踪的后续节点。核心取证：《清德宗实录》光绪二十八年相关条；宫中朱批奏折、内阁大库题本与《清史稿》相关志传；相关地区的档案、志书、契约、报刊或对方材料。该项锚定的是后续追踪问题，影响范围须按地区与群体分别判断。来源保留：战后、改革后或条约后的记忆、地方记录和官方总结常具有事后选择性；后续影响须按地点、群体和时间分层，不作整体化推断。
\eventanchor{EQ-1903-BANNER-POSTS-REFORM-PRELUDE}{「清廷废除部分旗人官职垄断，…}光绪二十九年（1903年），「清廷废除部分旗人官职垄断，调整旗籍制度与职务资格」之前的条件、信息和争议已通过中央与地方材料显现，构成该事件可追踪的前史节点。核心取证：《清德宗实录》光绪二十九年相关条；地方志、督抚奏折及地方档案；同时期地方档案、地方志、日记或对方材料应按具体地区补检。该项锚定的是前史条件与信息背景，不把条件直接写成唯一因果。来源保留：官修编年和地方材料的记录密度与立场并不相同；本行不将条件、传闻、呈报或后见解释直接等同为确定因果。
\eventanchor{EQ-1903-BANNER-POSTS-REFORM-DECISION}{围绕「清廷废除部分旗人官职垄…}光绪二十九年（1903年），围绕「清廷废除部分旗人官职垄断，调整旗籍制度与职务资格」的命令、调度、照会或呈报形成独立的中央—地方行动文书链。核心取证：《清德宗实录》光绪二十九年相关条；地方志、督抚奏折及地方档案。该项锚定的是命令或决策环节，命令抵达和结果仍须另外核验。来源保留：奏折、上谕、部议、军机档或条约可证明行政动作和机构立场，不自动证明所有地区、群体或对手按其意图行动。
\eventanchor{EQ-1903-BANNER-POSTS-REFORM-AFTERMATH}{「清廷废除部分旗人官职垄断，…}光绪二十九年（1903年），「清廷废除部分旗人官职垄断，调整旗籍制度与职务资格」引发的善后、执行或地方回应构成可由不同材料追踪的后续节点。核心取证：《清德宗实录》光绪二十九年相关条；地方志、督抚奏折及地方档案；相关地区的档案、志书、契约、报刊或对方材料。该项锚定的是后续追踪问题，影响范围须按地区与群体分别判断。来源保留：战后、改革后或条约后的记忆、地方记录和官方总结常具有事后选择性；后续影响须按地点、群体和时间分层，不作整体化推断。
\eventanchor{EQ-1903-SHANGHAI-COMMERCIAL-LAW-PRELUDE}{「商法、公司与工商管理的近代…}光绪二十九年（1903年），「商法、公司与工商管理的近代制度讨论在口岸和中央之间推进」之前的条件、信息和争议已通过中央与地方材料显现，构成该事件可追踪的前史节点。核心取证：《清德宗实录》光绪二十九年相关条；户部奏销、盐政、河工或海关档案；同时期地方档案、地方志、日记或对方材料应按具体地区补检。该项锚定的是前史条件与信息背景，不把条件直接写成唯一因果。来源保留：官修编年和地方材料的记录密度与立场并不相同；本行不将条件、传闻、呈报或后见解释直接等同为确定因果。
\eventanchor{EQ-1903-SHANGHAI-COMMERCIAL-LAW-DECISION}{围绕「商法、公司与工商管理的…}光绪二十九年（1903年），围绕「商法、公司与工商管理的近代制度讨论在口岸和中央之间推进」的命令、调度、照会或呈报形成独立的中央—地方行动文书链。核心取证：《清德宗实录》光绪二十九年相关条；户部奏销、盐政、河工或海关档案。该项锚定的是命令或决策环节，命令抵达和结果仍须另外核验。来源保留：奏折、上谕、部议、军机档或条约可证明行政动作和机构立场，不自动证明所有地区、群体或对手按其意图行动。
\eventanchor{EQ-1903-SHANGHAI-COMMERCIAL-LAW-AFTERMATH}{「商法、公司与工商管理的近代…}光绪二十九年（1903年），「商法、公司与工商管理的近代制度讨论在口岸和中央之间推进」引发的善后、执行或地方回应构成可由不同材料追踪的后续节点。核心取证：《清德宗实录》光绪二十九年相关条；户部奏销、盐政、河工或海关档案；相关地区的档案、志书、契约、报刊或对方材料。该项锚定的是后续追踪问题，影响范围须按地区与群体分别判断。来源保留：战后、改革后或条约后的记忆、地方记录和官方总结常具有事后选择性；后续影响须按地点、群体和时间分层，不作整体化推断。
\eventanchor{EQ-1904-RUSSO-JAPANESE-WAR-PRELUDE}{「日俄战争在中国东北爆发，清…}光绪三十年（1904年），「日俄战争在中国东北爆发，清廷名义中立而东北遭受战场化」之前的条件、信息和争议已通过中央与地方材料显现，构成该事件可追踪的前史节点。核心取证：《清德宗实录》光绪三十年相关条；军机处档、战事方略及地方奏报；同时期地方档案、地方志、日记或对方材料应按具体地区补检。该项锚定的是前史条件与信息背景，不把条件直接写成唯一因果。来源保留：官修编年和地方材料的记录密度与立场并不相同；本行不将条件、传闻、呈报或后见解释直接等同为确定因果。
\eventanchor{EQ-1904-RUSSO-JAPANESE-WAR-DECISION}{围绕「日俄战争在中国东北爆发…}光绪三十年（1904年），围绕「日俄战争在中国东北爆发，清廷名义中立而东北遭受战场化」的命令、调度、照会或呈报形成独立的中央—地方行动文书链。核心取证：《清德宗实录》光绪三十年相关条；军机处档、战事方略及地方奏报。该项锚定的是命令或决策环节，命令抵达和结果仍须另外核验。来源保留：奏折、上谕、部议、军机档或条约可证明行政动作和机构立场，不自动证明所有地区、群体或对手按其意图行动。
\eventanchor{EQ-1904-RUSSO-JAPANESE-WAR-AFTERMATH}{「日俄战争在中国东北爆发，清…}光绪三十年（1904年），「日俄战争在中国东北爆发，清廷名义中立而东北遭受战场化」引发的善后、执行或地方回应构成可由不同材料追踪的后续节点。核心取证：《清德宗实录》光绪三十年相关条；军机处档、战事方略及地方奏报；相关地区的档案、志书、契约、报刊或对方材料。该项锚定的是后续追踪问题，影响范围须按地区与群体分别判断。来源保留：战后、改革后或条约后的记忆、地方记录和官方总结常具有事后选择性；后续影响须按地点、群体和时间分层，不作整体化推断。
\eventanchor{EQ-1905-ABOLITION-OF-CIVIL-SERVICE-EXAM-PRELUDE}{「清廷废除科举，延续千余年的…}光绪三十一年（1905年），「清廷废除科举，延续千余年的取士制度终结」之前的条件、信息和争议已通过中央与地方材料显现，构成该事件可追踪的前史节点。核心取证：《清德宗实录》光绪三十一年相关条；内阁档、文集或报刊相关记录；同时期地方档案、地方志、日记或对方材料应按具体地区补检。该项锚定的是前史条件与信息背景，不把条件直接写成唯一因果。来源保留：官修编年和地方材料的记录密度与立场并不相同；本行不将条件、传闻、呈报或后见解释直接等同为确定因果。
\eventanchor{EQ-1905-ABOLITION-OF-CIVIL-SERVICE-EXAM-DECISION}{围绕「清廷废除科举，延续千余…}光绪三十一年（1905年），围绕「清廷废除科举，延续千余年的取士制度终结」的命令、调度、照会或呈报形成独立的中央—地方行动文书链。核心取证：《清德宗实录》光绪三十一年相关条；内阁档、文集或报刊相关记录。该项锚定的是命令或决策环节，命令抵达和结果仍须另外核验。来源保留：奏折、上谕、部议、军机档或条约可证明行政动作和机构立场，不自动证明所有地区、群体或对手按其意图行动。
\eventanchor{EQ-1905-ABOLITION-OF-CIVIL-SERVICE-EXAM-AFTERMATH}{「清廷废除科举，延续千余年的…}光绪三十一年（1905年），「清廷废除科举，延续千余年的取士制度终结」引发的善后、执行或地方回应构成可由不同材料追踪的后续节点。核心取证：《清德宗实录》光绪三十一年相关条；内阁档、文集或报刊相关记录；相关地区的档案、志书、契约、报刊或对方材料。该项锚定的是后续追踪问题，影响范围须按地区与群体分别判断。来源保留：战后、改革后或条约后的记忆、地方记录和官方总结常具有事后选择性；后续影响须按地点、群体和时间分层，不作整体化推断。
\eventanchor{EQ-1905-TONGMENGHUI-PRELUDE}{「中国同盟会在东京成立，革命…}光绪三十一年（1905年），「中国同盟会在东京成立，革命组织和宣传网络进一步整合」之前的条件、信息和争议已通过中央与地方材料显现，构成该事件可追踪的前史节点。核心取证：《清德宗实录》光绪三十一年相关条；地方志、督抚奏折及地方档案；同时期地方档案、地方志、日记或对方材料应按具体地区补检。该项锚定的是前史条件与信息背景，不把条件直接写成唯一因果。来源保留：官修编年和地方材料的记录密度与立场并不相同；本行不将条件、传闻、呈报或后见解释直接等同为确定因果。
\eventanchor{EQ-1905-TONGMENGHUI-DECISION}{围绕「中国同盟会在东京成立，…}光绪三十一年（1905年），围绕「中国同盟会在东京成立，革命组织和宣传网络进一步整合」的命令、调度、照会或呈报形成独立的中央—地方行动文书链。核心取证：《清德宗实录》光绪三十一年相关条；地方志、督抚奏折及地方档案。该项锚定的是命令或决策环节，命令抵达和结果仍须另外核验。来源保留：奏折、上谕、部议、军机档或条约可证明行政动作和机构立场，不自动证明所有地区、群体或对手按其意图行动。
\eventanchor{EQ-1905-TONGMENGHUI-AFTERMATH}{「中国同盟会在东京成立，革命…}光绪三十一年（1905年），「中国同盟会在东京成立，革命组织和宣传网络进一步整合」引发的善后、执行或地方回应构成可由不同材料追踪的后续节点。核心取证：《清德宗实录》光绪三十一年相关条；地方志、督抚奏折及地方档案；相关地区的档案、志书、契约、报刊或对方材料。该项锚定的是后续追踪问题，影响范围须按地区与群体分别判断。来源保留：战后、改革后或条约后的记忆、地方记录和官方总结常具有事后选择性；后续影响须按地点、群体和时间分层，不作整体化推断。
\eventanchor{EQ-1906-CONSTITUTIONAL-PROMISE-PRELUDE}{「清廷宣布预备立宪，着手考察…}光绪三十二年（1906年），「清廷宣布预备立宪，着手考察各国制度和改革官制」之前的条件、信息和争议已通过中央与地方材料显现，构成该事件可追踪的前史节点。核心取证：《清德宗实录》光绪三十二年相关条；宫中朱批奏折、内阁大库题本与《清史稿》相关志传；同时期地方档案、地方志、日记或对方材料应按具体地区补检。该项锚定的是前史条件与信息背景，不把条件直接写成唯一因果。来源保留：官修编年和地方材料的记录密度与立场并不相同；本行不将条件、传闻、呈报或后见解释直接等同为确定因果。
\eventanchor{EQ-1906-CONSTITUTIONAL-PROMISE-DECISION}{围绕「清廷宣布预备立宪，着手…}光绪三十二年（1906年），围绕「清廷宣布预备立宪，着手考察各国制度和改革官制」的命令、调度、照会或呈报形成独立的中央—地方行动文书链。核心取证：《清德宗实录》光绪三十二年相关条；宫中朱批奏折、内阁大库题本与《清史稿》相关志传。该项锚定的是命令或决策环节，命令抵达和结果仍须另外核验。来源保留：奏折、上谕、部议、军机档或条约可证明行政动作和机构立场，不自动证明所有地区、群体或对手按其意图行动。
\eventanchor{EQ-1906-CONSTITUTIONAL-PROMISE-AFTERMATH}{「清廷宣布预备立宪，着手考察…}光绪三十二年（1906年），「清廷宣布预备立宪，着手考察各国制度和改革官制」引发的善后、执行或地方回应构成可由不同材料追踪的后续节点。核心取证：《清德宗实录》光绪三十二年相关条；宫中朱批奏折、内阁大库题本与《清史稿》相关志传；相关地区的档案、志书、契约、报刊或对方材料。该项锚定的是后续追踪问题，影响范围须按地区与群体分别判断。来源保留：战后、改革后或条约后的记忆、地方记录和官方总结常具有事后选择性；后续影响须按地点、群体和时间分层，不作整体化推断。
\eventanchor{EQ-1907-MANCHURIAN-PROVINCES-PRELUDE}{「清廷将东三省调整为行省体制…}光绪三十三年（1907年），「清廷将东三省调整为行省体制，加强近代行政与边防经营」之前的条件、信息和争议已通过中央与地方材料显现，构成该事件可追踪的前史节点。核心取证：《清德宗实录》光绪三十三年相关条；军机处档、理藩院档或相关方略；同时期地方档案、地方志、日记或对方材料应按具体地区补检。该项锚定的是前史条件与信息背景，不把条件直接写成唯一因果。来源保留：官修编年和地方材料的记录密度与立场并不相同；本行不将条件、传闻、呈报或后见解释直接等同为确定因果。
\eventanchor{EQ-1907-MANCHURIAN-PROVINCES-DECISION}{围绕「清廷将东三省调整为行省…}光绪三十三年（1907年），围绕「清廷将东三省调整为行省体制，加强近代行政与边防经营」的命令、调度、照会或呈报形成独立的中央—地方行动文书链。核心取证：《清德宗实录》光绪三十三年相关条；军机处档、理藩院档或相关方略。该项锚定的是命令或决策环节，命令抵达和结果仍须另外核验。来源保留：奏折、上谕、部议、军机档或条约可证明行政动作和机构立场，不自动证明所有地区、群体或对手按其意图行动。
\eventanchor{EQ-1907-MANCHURIAN-PROVINCES-AFTERMATH}{「清廷将东三省调整为行省体制…}光绪三十三年（1907年），「清廷将东三省调整为行省体制，加强近代行政与边防经营」引发的善后、执行或地方回应构成可由不同材料追踪的后续节点。核心取证：《清德宗实录》光绪三十三年相关条；军机处档、理藩院档或相关方略；相关地区的档案、志书、契约、报刊或对方材料。该项锚定的是后续追踪问题，影响范围须按地区与群体分别判断。来源保留：战后、改革后或条约后的记忆、地方记录和官方总结常具有事后选择性；后续影响须按地点、群体和时间分层，不作整体化推断。
\eventanchor{EQ-1907-QING-LEGAL-REFORM-PRELUDE}{「清廷推进修律、审判和警政改…}光绪三十三年（1907年），「清廷推进修律、审判和警政改革，尝试改造传统法律制度」之前的条件、信息和争议已通过中央与地方材料显现，构成该事件可追踪的前史节点。核心取证：《清德宗实录》光绪三十三年相关条；宫中朱批奏折、内阁大库题本与《清史稿》相关志传；同时期地方档案、地方志、日记或对方材料应按具体地区补检。该项锚定的是前史条件与信息背景，不把条件直接写成唯一因果。来源保留：官修编年和地方材料的记录密度与立场并不相同；本行不将条件、传闻、呈报或后见解释直接等同为确定因果。
\eventanchor{EQ-1907-QING-LEGAL-REFORM-DECISION}{围绕「清廷推进修律、审判和警…}光绪三十三年（1907年），围绕「清廷推进修律、审判和警政改革，尝试改造传统法律制度」的命令、调度、照会或呈报形成独立的中央—地方行动文书链。核心取证：《清德宗实录》光绪三十三年相关条；宫中朱批奏折、内阁大库题本与《清史稿》相关志传。该项锚定的是命令或决策环节，命令抵达和结果仍须另外核验。来源保留：奏折、上谕、部议、军机档或条约可证明行政动作和机构立场，不自动证明所有地区、群体或对手按其意图行动。
\eventanchor{EQ-1907-QING-LEGAL-REFORM-AFTERMATH}{「清廷推进修律、审判和警政改…}光绪三十三年（1907年），「清廷推进修律、审判和警政改革，尝试改造传统法律制度」引发的善后、执行或地方回应构成可由不同材料追踪的后续节点。核心取证：《清德宗实录》光绪三十三年相关条；宫中朱批奏折、内阁大库题本与《清史稿》相关志传；相关地区的档案、志书、契约、报刊或对方材料。该项锚定的是后续追踪问题，影响范围须按地区与群体分别判断。来源保留：战后、改革后或条约后的记忆、地方记录和官方总结常具有事后选择性；后续影响须按地点、群体和时间分层，不作整体化推断。
\eventanchor{EQ-1908-GUANGXU-CIXI-DEATHS-PRELUDE}{「光绪帝与慈禧太后相继去世，…}光绪三十四年（1908年），「光绪帝与慈禧太后相继去世，溥仪继位为宣统帝」之前的条件、信息和争议已通过中央与地方材料显现，构成该事件可追踪的前史节点。核心取证：《清德宗实录》光绪三十四年相关条；宫中朱批奏折、内阁大库题本与《清史稿》相关志传；同时期地方档案、地方志、日记或对方材料应按具体地区补检。该项锚定的是前史条件与信息背景，不把条件直接写成唯一因果。来源保留：官修编年和地方材料的记录密度与立场并不相同；本行不将条件、传闻、呈报或后见解释直接等同为确定因果。
\eventanchor{EQ-1908-GUANGXU-CIXI-DEATHS-DECISION}{围绕「光绪帝与慈禧太后相继去…}光绪三十四年（1908年），围绕「光绪帝与慈禧太后相继去世，溥仪继位为宣统帝」的命令、调度、照会或呈报形成独立的中央—地方行动文书链。核心取证：《清德宗实录》光绪三十四年相关条；宫中朱批奏折、内阁大库题本与《清史稿》相关志传。该项锚定的是命令或决策环节，命令抵达和结果仍须另外核验。来源保留：奏折、上谕、部议、军机档或条约可证明行政动作和机构立场，不自动证明所有地区、群体或对手按其意图行动。
\eventanchor{EQ-1908-GUANGXU-CIXI-DEATHS-AFTERMATH}{「光绪帝与慈禧太后相继去世，…}光绪三十四年（1908年），「光绪帝与慈禧太后相继去世，溥仪继位为宣统帝」引发的善后、执行或地方回应构成可由不同材料追踪的后续节点。核心取证：《清德宗实录》光绪三十四年相关条；宫中朱批奏折、内阁大库题本与《清史稿》相关志传；相关地区的档案、志书、契约、报刊或对方材料。该项锚定的是后续追踪问题，影响范围须按地区与群体分别判断。来源保留：战后、改革后或条约后的记忆、地方记录和官方总结常具有事后选择性；后续影响须按地点、群体和时间分层，不作整体化推断。
\eventanchor{EQ-1908-CONSTITUTIONAL-OUTLINE-PRELUDE}{「清廷公布《钦定宪法大纲》等…}光绪三十四年（1908年），「清廷公布《钦定宪法大纲》等文件，明确预备立宪框架」之前的条件、信息和争议已通过中央与地方材料显现，构成该事件可追踪的前史节点。核心取证：《清德宗实录》光绪三十四年相关条；宫中朱批奏折、内阁大库题本与《清史稿》相关志传；同时期地方档案、地方志、日记或对方材料应按具体地区补检。该项锚定的是前史条件与信息背景，不把条件直接写成唯一因果。来源保留：官修编年和地方材料的记录密度与立场并不相同；本行不将条件、传闻、呈报或后见解释直接等同为确定因果。
\eventanchor{EQ-1908-CONSTITUTIONAL-OUTLINE-DECISION}{围绕「清廷公布《钦定宪法大纲…}光绪三十四年（1908年），围绕「清廷公布《钦定宪法大纲》等文件，明确预备立宪框架」的命令、调度、照会或呈报形成独立的中央—地方行动文书链。核心取证：《清德宗实录》光绪三十四年相关条；宫中朱批奏折、内阁大库题本与《清史稿》相关志传。该项锚定的是命令或决策环节，命令抵达和结果仍须另外核验。来源保留：奏折、上谕、部议、军机档或条约可证明行政动作和机构立场，不自动证明所有地区、群体或对手按其意图行动。
\eventanchor{EQ-1908-CONSTITUTIONAL-OUTLINE-AFTERMATH}{「清廷公布《钦定宪法大纲》等…}光绪三十四年（1908年），「清廷公布《钦定宪法大纲》等文件，明确预备立宪框架」引发的善后、执行或地方回应构成可由不同材料追踪的后续节点。核心取证：《清德宗实录》光绪三十四年相关条；宫中朱批奏折、内阁大库题本与《清史稿》相关志传；相关地区的档案、志书、契约、报刊或对方材料。该项锚定的是后续追踪问题，影响范围须按地区与群体分别判断。来源保留：战后、改革后或条约后的记忆、地方记录和官方总结常具有事后选择性；后续影响须按地点、群体和时间分层，不作整体化推断。
